\begin{enumerate}
  \item[\textbf{Q1.}] The paper builds $U_m$ by an $m$-deep recursion whose
    gate cost is $3^{m}$ oracle queries but whose \emph{ancilla + phase-shift}
    cost is not given explicitly. On real Clifford+T hardware, is the
    $T$-count of one recursion step of $U R_s U^\dagger R_t U$ larger,
    smaller, or comparable to the $T$-count of the equivalent
    $q_m$-query standard-Grover circuit that would give the same
    success probability with an ideal iteration count? A concrete
    counting for $m=1,\dots,4$ on a fault-tolerant compiler would
    settle this.
  \item[\textbf{Q2.}] With realistic depolarising or coherent gate
    noise per Clifford of $\sim 10^{-3}$, at what recursion depth $m^{*}$
    does the accumulated gate error \emph{overtake} the algorithmic
    error $\epsilon^{\,3^{m}}$, so that further recursion \emph{hurts}?
    The paper's fixed-point property is a statement in the noiseless
    unitary limit; the practical crossover $m^{*}(p_\text{gate},\epsilon)$
    is a natural follow-on our reproduction did not measure.
  \item[\textbf{Q3.}] Grover claims (Sec.~5) that averaging over the
    $f \in [0.75, 1.0]$ uniform prior yields a failure-probability
    improvement of ``approximately one fourth'' relative to the best
    known quantum single-query scheme (Younes--Rowe--Miller~2003). Our
    analytic integration reproduces the qualitative $\sim 4\times$
    factor but the exact averaging convention that turns $(1-f)^3$
    into ``$\approx 0.8\%$'' is not spelled out in the paper. What is
    the precise conditional-on-success sampling convention Grover is
    using, and does his stated $0.8\%$ vs. $3.12\%$ hold under a
    strictly uniform-$f$ marginalisation?
  \item[\textbf{Q4.}] The recursion generates one distinct unitary $U_m$
    per depth $m$ (unlike standard amplitude amplification, which
    reuses a single $Q$). We observed that the fixed-point property
    depends on this asymmetry between $R_s$ and $R_s^\dagger$ in the
    expanded product. If one \emph{approximates} $R_s^\dagger \approx R_s$
    (i.e.\ uses only $R_s$ throughout, dropping the daggers on the
    inner shifts), how much does the fixed-point property degrade as a
    function of $m$? Numerically bounding this would clarify how
    strict the ``different transformation each step'' requirement is.
  \item[\textbf{Q5.}] The paper (Sec.~6) promises that the same
    transformation reduces \emph{systematic} error in an arbitrary
    driving unitary $U$ from $\epsilon$ to $\epsilon^{3}$, provided
    $U^\dagger$ can be implemented with the same error. Does this
    fixed-point advantage survive when $U^\dagger$ is only known to
    the same imprecision as $U$ (i.e.\ when the daggered call carries
    an independent $O(\epsilon)$ error rather than being the exact
    inverse)? A short numerical experiment perturbing $U$ and
    $U^\dagger$ independently would resolve whether Section~6's
    error-correction promise is robust or fragile in the realistic
    ``mis-calibrated inverse'' regime.
\end{enumerate}
